% ===== PRÉAMBULE CANONIQUE — à coller VERBATIM en tête de CHAQUE .tex =====
\documentclass[11pt,a4paper]{article}
\usepackage[utf8]{inputenc}\usepackage[T1]{fontenc}\usepackage{lmodern}
\usepackage[french]{babel}
\usepackage{amsmath,amssymb,mathtools}\usepackage{icomma}
\usepackage[version=4]{mhchem}          % \ce{...} : équations & formules
\usepackage{siunitx}\sisetup{locale=FR} % \qty{}{}, \si{}, \ang{}
\usepackage{chemfig}                    % \chemfig{...} : structures développées
\usepackage{xcolor}\usepackage{colortbl}
\usepackage{tikz}\usepackage{pgfplots}\pgfplotsset{compat=1.18}
\usepackage[most]{tcolorbox}\usepackage{enumitem}\usepackage{array,booktabs}
\usepackage[a4paper,margin=2.2cm]{geometry}
\usetikzlibrary{arrows.meta,calc,angles,quotes,positioning,patterns,backgrounds,shapes.geometric}
\definecolor{primary}{RGB}{222,49,99}\definecolor{primarydark}{RGB}{150,22,55}
\definecolor{defcol}{RGB}{0,114,178}\definecolor{methcol}{RGB}{52,92,148}
\definecolor{excol}{RGB}{0,135,90}\definecolor{attncol}{RGB}{200,40,55}
\tcbset{boxbase/.style={enhanced,breakable,boxrule=0.9pt,arc=2.5pt,
  left=3mm,right=3mm,top=2.3mm,bottom=2.3mm,fonttitle=\bfseries,coltitle=white}}
% --- boîtes TUTORIEL (noms EXACTS, ne pas renommer) ---
\newtcolorbox{cle}[1][]{boxbase,colback=defcol!6,colframe=defcol,
  title={\textsc{À retenir}\ifx\relax#1\relax\relax\else\textmd{ : \itshape #1}\fi}}
\newtcolorbox{methode}[1][]{boxbase,colback=methcol!7,colframe=methcol,sharp corners,
  title={\textsc{Méthode}\ifx\relax#1\relax\relax\else\textmd{ --- \itshape #1}\fi}}
\newtcolorbox{exemplebox}[1][]{boxbase,colback=excol!5,colframe=excol,
  title={\textsc{Exemple}\ifx\relax#1\relax\relax\else\textmd{ : \itshape #1}\fi}}
\newtcolorbox{piege}[1][]{boxbase,colback=attncol!6,colframe=attncol,
  title={\textsc{Piège}\ifx\relax#1\relax\relax\else\textmd{ : \itshape #1}\fi}}
\newtcolorbox{remarque}[1][]{boxbase,colback=black!4,colframe=black!45,
  title={\textsc{Remarque}\ifx\relax#1\relax\relax\else\textmd{ : \itshape #1}\fi}}
% --- boîtes EXERCICE / CORRIGÉ (noms EXACTS, auto-comptées) ---
\newtcolorbox[auto counter]{exo}[1][]{boxbase,colback=primary!4,colframe=primary,
  title={\textsc{Exercice}~\thetcbcounter\ifx\relax#1\relax\relax\else\textmd{ --- \itshape #1}\fi}}
\newtcolorbox[auto counter]{corrige}[1][]{boxbase,colback=excol!4,colframe=excol,
  title={\textsc{Corrigé}~\thetcbcounter\ifx\relax#1\relax\relax\else\textmd{ --- \itshape #1}\fi}}
\newcommand{\R}{\mathbb{R}}\newcommand{\N}{\mathbb{N}}\newcommand{\Z}{\mathbb{Z}}\newcommand{\ds}{\displaystyle}
% =========================================================================
\begin{document}
\sisetup{per-mode=symbol}
\begin{exo}[la bonne expression parmi quatre]
On propose quatre expressions candidates pour le produit de solubilité de \ce{Ag3PO4} :
\begin{center}
\begin{tabular}{ll ll}
\textbf{A.} $3\,[\ce{Ag+}]\,[\ce{PO4^{3-}}]$ &\quad&
\textbf{B.} $[\ce{Ag^{3+}}]\,[\ce{PO4^{3-}}]$ \\[4pt]
\textbf{C.} $[\ce{Ag+}]^{3}+[\ce{PO4^{3-}}]$ &\quad&
\textbf{D.} $[\ce{Ag+}]^{3}\,[\ce{PO4^{3-}}]$ \\
\end{tabular}
\end{center}
\begin{enumerate}[label=\alph*)]
  \item Écris l'équation de dissolution de \ce{Ag3PO4}.
  \item Indique la \textbf{seule} expression correcte.
  \item Pour chacune des trois autres, explique précisément l'erreur.
\end{enumerate}
\end{exo}

\begin{exo}[assertions sur les unités : \ce{Pb3(PO4)2}]
On considère le phosphate de plomb \ce{Pb3(PO4)2}. Pour chaque assertion, dis si elle est correcte et
justifie.
\begin{enumerate}[label=\alph*)]
  \item L'unité de la solubilité de \ce{Pb3(PO4)2} est $\mathrm{M}^{5}$.
  \item L'unité du produit de solubilité de \ce{Pb3(PO4)2} est $\mathrm{mol^{5}\,L^{-5}}$.
  \item Le produit de solubilité est une constante d'équilibre.
  \item $K_{\mathrm{ps}}$ varie avec la température.
\end{enumerate}
\end{exo}

\begin{exo}[trois solutions saturées analysées]
Un laboratoire mesure les concentrations ioniques dans trois solutions saturées. Calcule le
$K_{\mathrm{ps}}$ de chaque sel.
\begin{enumerate}[label=\alph*)]
  \item \ce{CaF2} : $[\ce{Ca^{2+}}]=\qty{2.0e-4}{\mole\per\litre}$, $[\ce{F-}]=\qty{4.0e-4}{\mole\per\litre}$.
  \item \ce{Ag3PO4} : $[\ce{Ag+}]=\qty{6.0e-5}{\mole\per\litre}$, $[\ce{PO4^{3-}}]=\qty{2.0e-5}{\mole\per\litre}$.
  \item \ce{Ca3(PO4)2} : $[\ce{Ca^{2+}}]=\qty{3.0e-7}{\mole\per\litre}$, $[\ce{PO4^{3-}}]=\qty{2.0e-7}{\mole\per\litre}$.
\end{enumerate}
\end{exo}

\begin{exo}[remonter à la formule du sel]
On te donne l'expression du produit de solubilité d'un sel inconnu :
\[ K_{\mathrm{ps}} = [\ce{M^{3+}}]^{2}\,[\ce{SO4^{2-}}]^{3}. \]
\begin{enumerate}[label=\alph*)]
  \item Quelle est la formule brute du sel ?
  \item Écris son équation de dissolution.
  \item Ce sel est-il un sel de type $\mathrm{MX}_n$ (un seul cation) ? Quel est son type $p{:}q$ ?
\end{enumerate}
\end{exo}
\end{document}
